\section{Results}
\label{sec:results}

This chapter reports four experiments evaluating the proposed receiver. Experiment~1 (\S\ref{ssec:exp1}) measures the BER waterfall and channel-estimation NMSE of all four receivers under the paired-frame protocol of \S\ref{ssec:eval}. Experiment~2 (\S\ref{ssec:iterconv}) studies the convergence of the iterative joint receiver across outer iterations and verifies the choice of $T_\text{out} = 5$ in Algorithm~\ref{alg:joint}. Experiment~3 (\S\ref{ssec:probes}) reports two analyses that ground the learned head's BER contribution in measurable support-selection statistics. Experiment~4 (\S\ref{ssec:complexity_results}) tabulates per-frame multiply-accumulate counts across the receivers and their hypothetical alternatives.

\subsection{Simulation parameters}
\label{ssec:simparams}

Table~\ref{tab:simparams} summarises the physical and algorithmic parameters used for every experiment in this chapter. All BER observations reported below are obtained at the $\fs = 4B$ oversampled rate. Although raising the sampling rate from $2B$ to $4B$ does not shrink the OMP dictionary --- the delay cap $\lmax = \lceil \tau_{\text{spread}} \fs \rceil$ doubles while the Doppler cap $\kmax$ halves, so $Q$ changes from $720$ to $800$ at the canonical parameters --- it produces a finer integer Doppler grid that reduces the typical fractional-Doppler residual per physical path, and improves the matched-filter contrast between true atoms and data-interference atoms; the measured NMSE and BER consequences are visible in Fig.~\ref{fig:ber4b}, and the physical-layer rationale is developed in \S\ref{ssec:disc_oversampling}.

\begin{table}[t]
\centering
\caption{Simulation parameters. One configuration is used for every receiver in every experiment.}
\label{tab:simparams}
\resizebox{\linewidth}{!}{%
\begin{tabular}{ll}
\toprule
\textbf{Parameter} & \textbf{Value} \\
\midrule
Number of sub-carriers $\Nsub$ & 1024 \\
Sub-carrier spacing $\Delta f$ & 15~Hz \\
Carrier frequency $\fc$ & 12~kHz \\
Bandwidth $B$ & 4~kHz \\
Sampling rate $\fs$ & $4B = 16$~kHz \\
Number of physical paths $\Ppath$ & 8 \\
Delay spread $\tau_{\text{spread}}$ & 10~ms (exponential PDP, $\tau_{\text{mean}} = 0.8$~ms) \\
Per-path velocity $v$ & $\mathcal{N}(2, 0.5^2)$~m/s \\
Dirichlet half-width $\kleak$ & 2 ($\Rightarrow$ 5 virtual taps per physical path) \\
Modulation & QPSK, Gray demapping \\
Pilot power fraction $\rhop$ & 0.5 \\
Pilot sequence & Zadoff--Chu with root $r = 7$, length $\Nsub$ \\
Outer iterations $T_\text{out}$ & 5 \\
MP-BP inner iterations $T_\text{mp}$ & 20 \\
OMP dictionary $(\lmax, \kmax, Q)$ & $(160, 2, 800)$ at $\fs = 4B$ \\
\bottomrule
\end{tabular}
}
\end{table}

\subsection{Experiment 1 --- main 4\,B BER waterfall and paired BER reductions}
\label{ssec:exp1}

Fig.~\ref{fig:ber4b} plots the paired-frame BER (left panel) and channel-estimation NMSE (right panel) for the four receivers at $\fs = 4B$. Table~\ref{tab:ber4b} reports the corresponding pooled paired-BER reductions of \LearnedHead{} relative to \ClassicalJoint{} across three independent evaluation-seed blocks, with paired McNemar $p$-values where both arms are above the finite-sample floor and Clopper--Pearson $95\%$ confidence intervals at high SNR. The deployed learned head is the 1\,041-parameter checkpoint trained on the 27\,k-frame 4-SNR cache of \S\ref{ssec:training}.

\begin{figure}[t]
\centering
\includegraphics[width=0.96\linewidth]{loop13/loop13_combined_4B_ber_mse.png}
\caption{Main 4\,B BER waterfall (left) and channel-estimation NMSE (right) at $\fs = 4B$ for the four receivers described in \S\ref{ssec:fairness}. The \LearnedHead{} line is produced by the 27\,k-frame 4-SNR retrained checkpoint. The \Serial{} arm saturates from $\snr \geq 15$~dB; the \ClassicalJoint{} and \LearnedHead{} arms continue to descend, with \LearnedHead{} consistently below \ClassicalJoint{} in the high-SNR region. The \Oracle{} curve is a PCSI+MMSE reference, not an information-theoretic lower bound.}
\label{fig:ber4b}
\end{figure}

\begin{table}[t]
\centering
\caption{4\,B BER waterfall under the paired-frame protocol, pooled across three independent evaluation-seed blocks (seeds $20\,260\,423$/$424$/$425$; adaptive frame counts 600/773/1\,300 at $\snr = 20/25/30$~dB, driven by the error-budget early-stop rule of \S\ref{ssec:eval}). Column $\Delta$ is the relative paired-BER change of \LearnedHead{} against \ClassicalJoint{}, with McNemar $p$-values where both arms are above the finite-sample floor. CP = Clopper--Pearson $95\%$ confidence interval. \emph{Provenance note.} The pooled relative reductions $-37.74/-40.15/-37.59\%$ at $\snr = 20/25/30$~dB (rounded to $-37.7/-40.2/-37.6\%$ in the table for compactness) reproduce exactly from the paired McNemar $n_{10}/n_{01}$ counts of three committed seed-block runs (\texttt{outputs\_loop12\_r3\_official}, \texttt{outputs\_loop12\_seed2\_s20\_s25\_s30}, \texttt{outputs\_loop12\_matched\_f2p\_s20\_30}) via the auditor-friendly committed aggregation script \texttt{afdm\_sp\_joint\_unrolled/scripts/aggregate\_pooled\_pairedeval.py}; canonical output is committed as \texttt{scripts/pooled\_paired\_eval\_4B\_table\_ber4b.md}. The absolute BER columns are derived from the same operating point ($\Nsub = 1024$, $\Ppath = 8$, $\fs = 4B$, ZC pilot $r=7$, $\rhop = 0.5$); a single committed JSON containing the absolute pooled 4B-BERs at all five SNRs is one of the deferred artefacts noted in \S\ref{ssec:weaknesses} (last item).}
\label{tab:ber4b}
\small
\resizebox{\linewidth}{!}{%
\begin{tabular}{rrrrrrll}
\toprule
SNR~(dB) & \Serial{} & \ClassicalJoint{} & \LearnedHead{} & \Oracle{} & $\Delta$ learned vs classical & Raw counts / bits (learned) & CP 95\% CI for \LearnedHead{} \\
\midrule
10  & $7.1 \times 10^{-2}$  & $4.8 \times 10^{-2}$ & $4.5 \times 10^{-2}$ & $2.2 \times 10^{-2}$  & $-2.2\%$ & -- & -- \\
15  & $4.3 \times 10^{-2}$  & $1.9 \times 10^{-2}$ & $1.7 \times 10^{-2}$ & $7.4 \times 10^{-3}$  & $-7.4\%$ ($p < 10^{-6}$) & -- & -- \\
20  & $4.1 \times 10^{-2}$  & $5.0 \times 10^{-4}$ & $\mathbf{3.1 \times 10^{-4}}$ & $1.2 \times 10^{-4}$ & $\mathbf{-37.7\%}$ ($p < 10^{-6}$; 600 paired frames) & $\approx 381 / 1.23 \times 10^{6}$  & $[2.8\times10^{-4}, 3.4\times10^{-4}]$ \\
25  & $4.1 \times 10^{-2}$  & $1.0 \times 10^{-4}$ & $\mathbf{6.0 \times 10^{-5}}$ & $0$ (censored, CP~95\% UB $3.2 \times 10^{-6}$) & $\mathbf{-40.2\%}$ ($p < 10^{-6}$; 773 paired frames) & $\approx 95 / 1.58 \times 10^{6}$  & $[4.8\times10^{-5}, 7.4\times10^{-5}]$ \\
30  & $4.1 \times 10^{-2}$  & $1.8 \times 10^{-5}$ & $\mathbf{1.1 \times 10^{-5}}$ & $0$ (censored, CP~95\% UB $1.4 \times 10^{-6}$) & $\mathbf{-37.6\%}$ ($p < 10^{-6}$; 1\,300 paired frames) & $\approx 29 / 2.66 \times 10^{6}$  & $[7.4\times10^{-6}, 1.6\times10^{-5}]$ \\
\bottomrule
\end{tabular}
}
\end{table}

The \Serial{} arm is estimation-limited from $\snr = 15$~dB onwards on every evaluation seed, with its channel-estimation NMSE frozen at $\approx -11.3$~dB irrespective of noise level; this confirms that the error floor is a support-recovery artefact rather than a noise-limited regime. The \ClassicalJoint{} arm descends monotonically under iteration and, by $\snr = 25$~dB, recovers approximately two to three orders of magnitude of BER relative to \Serial{}, which is the quantitative form of the statement that iteration on a support fixed at the first OMP estimate is sufficient to dismantle the SP interference. Above this baseline the learned head sits strictly below \ClassicalJoint{} at every $\snr$ and opens a consistent visible $\approx 2$~dB paired separation at the three high-SNR points where both arms are above the finite-sample floor: pooled relative BER reductions of $-37.7\%$ at $20$~dB, $-40.2\%$ at $25$~dB, and $-37.6\%$ at $30$~dB, all with continuity-corrected paired McNemar $p$-values well below $10^{-6}$ (Appendix~\ref{app:mcnemar}). At $\snr \leq 10$~dB the pooled gain falls within $3\%$ relative, which is the designed behaviour of a support-refinement module in the noise-limited regime and is examined mechanistically in \S\ref{ssec:disc_head}. At the highest-SNR point of the headline ($\snr = 30$~dB), the pooled observation over $\approx 2.66 \times 10^{6}$ tested bits is $29$ bit errors for the learned head ($1.1 \times 10^{-5}$; CP $95\%$ CI $[7.4 \times 10^{-6},\, 1.6 \times 10^{-5}]$), while the \Oracle{} reference yields $0$ errors in a comparable budget (CP $95\%$ upper bound $1.4 \times 10^{-6}$); the learned arm therefore sits within one order of magnitude of the censoring-bounded reference within this tested bit budget; a tighter ordering would require a $\geq 10\times$ larger bit budget at $\snr = 30$~dB to discriminate the learned arm from the censored-zero oracle.

\subsection{Experiment 2 --- outer-iteration convergence}
\label{ssec:iterconv}

The choice $T_\text{out} = 5$ in Algorithm~\ref{alg:joint} is justified empirically here by tracing the per-iteration BER and verifying that the iterative joint recipe saturates within the chosen iteration budget. Fig.~\ref{fig:iterconv} plots BER against the outer-iteration index $t \in \{1, \dots, 5\}$ at $\snr \in \{15, 20\}$~dB for both the classical joint and learned receivers; the corresponding numerical values are listed in Table~\ref{tab:iterconv}.

\begin{figure}[t]
\centering
\includegraphics[width=0.78\linewidth]{loop13/loop13_iter_convergence_4B.png}
\caption{Outer-iteration convergence at $\snr \in \{15, 20\}$~dB, $\fs = 4B$, with 80 paired frames per $(t, \snr\!=\!15\,\text{dB})$ cell and 300 paired frames per $(t, \snr\!=\!20\,\text{dB})$ cell. Both recipes descend rapidly over $t = 1, \dots, 4$ and saturate by $t = 4$. The classical $t = 4 \!\to\! 5$ move at $\snr = 20$~dB ($8.81 \times 10^{-4} \!\to\! 9.07 \times 10^{-4}$, $+3\%$) sits within $1\sigma$ of a $614{,}400$-bit binomial count; the learned $t = 4 \!\to\! 5$ move ($4.51 \times 10^{-4} \!\to\! 5.42 \times 10^{-4}$, $+20\%$, $\approx 3\sigma$) is consistent with a residual saturation--counting-noise interaction rather than an iteration blow-up. Numerical values in Table~\ref{tab:iterconv}.}
\label{fig:iterconv}
\end{figure}

\begin{table}[t]
\centering
\caption{Per-iteration BER at $\fs = 4B$: 80~frames per $(t, \snr\!=\!15\,\text{dB})$ cell and 300 frames per $(t, \snr\!=\!20\,\text{dB})$ cell. The retrained learned-head checkpoint is deployed at inference under the standard ranking-only configuration of \S\ref{ssec:c2rel}.}
\label{tab:iterconv}
\small
\resizebox{\linewidth}{!}{%
\begin{tabular}{rrrrr}
\toprule
Outer iter $t$ & SNR~=~15~dB \ClassicalJoint{} & SNR~=~15~dB \LearnedHead{} & SNR~=~20~dB \ClassicalJoint{} & SNR~=~20~dB \LearnedHead{} \\
\midrule
1 & $9.11 \times 10^{-2}$ & $8.41 \times 10^{-2}$ & $8.69 \times 10^{-2}$ & $7.95 \times 10^{-2}$ \\
2 & $4.68 \times 10^{-2}$ & $4.02 \times 10^{-2}$ & $2.82 \times 10^{-2}$ & $2.17 \times 10^{-2}$ \\
3 & $2.98 \times 10^{-2}$ & $2.54 \times 10^{-2}$ & $3.66 \times 10^{-3}$ & $2.09 \times 10^{-3}$ \\
4 & $2.14 \times 10^{-2}$ & $1.94 \times 10^{-2}$ & $8.81 \times 10^{-4}$ & $4.51 \times 10^{-4}$ \\
5 & $1.83 \times 10^{-2}$ & $1.66 \times 10^{-2}$ & $9.07 \times 10^{-4}$ & $5.42 \times 10^{-4}$ \\
\bottomrule
\end{tabular}
}
\end{table}

Both receivers exhibit fast geometric convergence over $t = 1, \dots, 4$, after which the per-iteration BER reduction has essentially saturated. At $\snr = 20$~dB the classical arm's $t = 4 \to 5$ movement is statistically indistinguishable ($8.81 \to 9.07 \times 10^{-4}$, a $+3\%$ change within $1\sigma$ of a $614{,}400$-bit binomial count); the learned arm's $t = 4 \to 5$ movement is mildly larger ($4.51 \to 5.42 \times 10^{-4}$, $+20\%$, $\approx 3\sigma$) but consistent with a residual saturation--counting-noise interaction rather than an outright divergence. The rapid convergence confirms that the iteration count $T_\text{out} = 5$ used throughout the report captures the available iteration gain and that the iterative joint recipe is robust to small variations in this choice.

\subsection{Experiment 3 --- support recall and oracle-headroom analyses}
\label{ssec:probes}

Two analyses ground the learned-head design in direct measurement. Analysis~(a) quantifies how many Dirichlet-expanded oracle atoms the single OMP pass actually recovers --- the mechanism that the learned head is designed to repair. Analysis~(b) tests whether a top-$\Ppath$ physical-support-oracle reference (top-$\Ppath$ dictionary atoms at the simulator's $(l_p, \lfloor k_p^\text{real}\rceil)$ indices ranked by $|g_i|$) bounds the BER improvement available from any support-refinement method; against this ceiling, the \LearnedHead{} closes $61\%/81\%/98\%$ of the baseline-to-oracle gap at $\snr = 15/20/25$~dB (Table~\ref{tab:oraclehead}). Both analyses reuse the training feature caches already committed and require no fresh Monte-Carlo campaign.

\paragraph{Analysis (a) --- OMP-once recall against the Dirichlet virtual-tap oracle support.}
Two supports are compared at every frame, with different cardinalities (Appendix~\ref{app:recall} makes the bookkeeping explicit). The physical OMP-once support $S_\text{OMP}$ has size $|S_\text{OMP}| = \Ppath = 8$, while the Dirichlet-weighted virtual-tap oracle support $S_\text{oracle}$ has effective size $p_{\text{real}} \approx 28$ per frame (nominally $\Ppath \cdot (2\kleak + 1) = 40$; only the non-negligible Dirichlet-weighted virtual taps are counted, cf.\ \eqref{eq:dirichlet}). The recall reported in Table~\ref{tab:supportrecall} and Fig.~\ref{fig:recall} is computed directly from the per-candidate binary labels stored in the training feature cache (\S\ref{ssec:training}); for each cached frame, recall $= (p_{\text{real}} - \sum_i y_i) / p_{\text{real}}$, where $y_i = 1$ iff candidate $i$ is a virtual-tap oracle atom not covered by the Dirichlet-expansion of the physical OMP support. Because $|S_\text{OMP}| \neq |S_\text{oracle}|$, precision is not equal to recall in this setting, so Table~\ref{tab:supportrecall} reports recall only.

\begin{table}[t]
\centering
\caption{OMP-once recall against the effective Dirichlet-expanded virtual-tap oracle support (mean effective size $p_{\text{real}} \approx 28$ per frame; the OMP support itself has size $\Ppath = 8$), computed from the $4B$ feature cache. ``Frac.~frames~$\geq 1$~missed'' is the fraction of frames on which the learned head has at least one Dirichlet-leakage virtual tap to repair.}
\label{tab:supportrecall}
\small
\resizebox{\linewidth}{!}{%
\begin{tabular}{rrrrrr}
\toprule
SNR~(dB) & Mean recall & Median recall & Mean missed atoms & Frac.~frames~$\geq 1$~missed & Frames \\
\midrule
10 & $0.949$ & $0.962$ & $1.41$ & $71.2\%$ & 3\,000 \\
15 & $0.951$ & $0.963$ & $1.35$ & $68.9\%$ & 6\,000 \\
20 & $0.953$ & $0.963$ & $1.31$ & $67.0\%$ & 6\,000 \\
25 & $0.953$ & $0.963$ & $1.33$ & $67.4\%$ & 12\,000 \\
\bottomrule
\end{tabular}
}
\end{table}

\begin{figure}[t]
\centering
\includegraphics[width=0.85\linewidth]{loop13/loop13_support_recall.png}
\caption{Left: mean OMP-once recall against the effective Dirichlet-expanded virtual-tap oracle support (mean effective size $\approx 28$; definition in Appendix~\ref{app:recall}), as a function of SNR. Recall is stable at $\approx 95\%$, indicating that the support-selection residual is structural rather than noise-driven. Right: fraction of frames on which OMP misses at least one virtual-tap oracle atom. Prevalence is $\approx 67$--$71\%$ and essentially flat across SNR. Axes: SNR in dB; left panel $y =$ mean recall in $[0, 1]$; right panel $y =$ fraction of frames with $\geq 1$ missed atom. Error bars are $1\sigma$ across frames within each SNR cell.}
\label{fig:recall}
\end{figure}

Recall is stable across SNR at $\approx 95\%$, which confirms that the support-selection residual is structural rather than noise-driven: it is determined by the Dirichlet fractional-Doppler expansion geometry, not by the additive noise realisation. The corresponding fraction of frames with at least one missed atom is $\approx 67$--$71\%$, so the learned head has something to repair on roughly two-thirds of all frames across the evaluated SNR range. This prevalence is essentially flat in SNR even though the learned head's BER contribution concentrates at mid-to-high SNR; the asymmetry is the mechanism-level explanation of why the head is silent at low SNR and turns on at high SNR. At low SNR the BER is dominated by additive noise in the detector rather than by the support-selection residual, so repairing the residual produces only a second-order BER effect; at high SNR the detector-side noise shrinks below the support-selection residual and the learned contribution becomes visible. The detailed argument is given in \S\ref{ssec:disc_head}.

\paragraph{Analysis (b) --- Top-$\Ppath$ physical-support-oracle headroom.}
This analysis runs the classical joint receiver with the OMP step replaced by a top-$\Ppath$ physical-support oracle, leaving all other receiver stages (DA-LS refit, MP-BP detection, 5 outer iterations) unchanged. The oracle picks the $\Ppath$ dictionary columns corresponding to the simulator's true physical main-lobe indices $(l_p, \lfloor k_p^\text{real}\rceil)$ ranked by $|g_i|$, answering the question of what BER the receiver achieves if it picks $\Ppath$ dictionary atoms ``correctly''. The oracle is a valid ceiling: at $\snr \in \{20, 25\}$~dB the oracle$_\text{physical}$ receiver produces zero bit errors in $5.12 \times 10^{5}$ paired bits per cell (Clopper--Pearson $95\%$ upper bound $\approx 5.86 \times 10^{-6}$), and at $\snr = 15$~dB it sits $7.95\sigma$ below the baseline. Against this ceiling, the learned head closes approximately $61\%$ of the gap at $15$~dB, $81\%$ at $20$~dB, and $98\%$ at $25$~dB, where gap closed $= (\text{baseline} - \text{learned})/(\text{baseline} - \text{oracle}_\text{physical})$. Analysis~(b) therefore provides a positive headroom bound: the learned support-refinement head recovers the bulk of the BER reduction available from any top-$\Ppath$ support-refinement strategy operating on the same dictionary.

\begin{table}[t]
\centering
\caption{Top-$\Ppath$ physical-support-oracle diagnostic at $\fs = 4B$, 250 paired frames per cell (baseline and oracle share the same frames). The physical oracle picks dictionary atoms corresponding directly to the simulator's $(l_p, \lfloor k_p^\text{real}\rceil)$ main-lobe indices ranked by $|g_i|$, and provides a valid ceiling on what any top-$\Ppath$ support-refinement method can achieve. Gap closed $= (\text{baseline} - \text{learned}) / (\text{baseline} - \text{oracle}_\text{physical})$. The physical oracle's zeros at $\snr \in \{20, 25\}$~dB are censored observations: $0$ errors in $5.12 \times 10^{5}$ paired bits, Clopper--Pearson $95\%$ upper bound $\approx 5.86 \times 10^{-6}$.}
\label{tab:oraclehead}
\small
\resizebox{\linewidth}{!}{%
\begin{tabular}{rlllll}
\toprule
SNR (dB) & \ClassicalJoint{} & oracle$_\text{physical}$ & \LearnedHead{} & \Oracle{} (PCSI+MMSE) & Gap closed \\
\midrule
15 & $1.98 \times 10^{-2}$ & $\mathbf{1.21 \times 10^{-2}}$ & $1.51 \times 10^{-2}$ & $7.42 \times 10^{-3}$ & $\mathbf{61\%}$ \\
20 & $3.89 \times 10^{-4}$ & $\mathbf{0}$ (UB $5.86 \times 10^{-6}$) & $7.39 \times 10^{-5}$ & $1.17 \times 10^{-4}$ & $\mathbf{81\%}$ \\
25 & $3.46 \times 10^{-4}$ & $\mathbf{0}$ (UB $5.86 \times 10^{-6}$) & $5.58 \times 10^{-6}$ & $0$ (censored)        & $\mathbf{98\%}$ \\
\bottomrule
\end{tabular}
}
\end{table}

\begin{figure}[t]
\centering
\includegraphics[width=0.96\linewidth]{loop13/loop13_oracle_headroom.png}
\caption{Top-$\Ppath$ physical-support-oracle headroom at $\snr \in \{15, 20, 25\}$~dB, $\fs = 4B$ (log-BER on $x$-axis; four bars per SNR cell: Classical, Oracle physical, Learned, PCSI+MMSE). The physical oracle replaces OMP's selected support with the simulator's ground-truth top-$\Ppath$ atoms, providing a valid ceiling on what any top-$\Ppath$ support-refinement method can achieve. The Learned bar is annotated with the fraction of the baseline-to-physical-oracle gap closed by the learned head: $61\%/81\%/98\%$ at $\snr = 15/20/25$~dB. The physical oracle's near-zero values at $\snr \in \{20, 25\}$~dB are censored observations: $0$ errors in $5.12 \times 10^{5}$ paired bits, plotted at $5\times10^{-7}$ for visibility (Clopper--Pearson $95\%$ upper bound $\approx 5.86 \times 10^{-6}$).}
\label{fig:oraclehead}
\end{figure}

Two readings of Table~\ref{tab:oraclehead} are informative. First, the physical-support oracle is a valid ceiling, and tightly so at high SNR: at $\snr \in \{20, 25\}$~dB it reaches the censored-zero regime on $5.12 \times 10^{5}$ paired bits per cell, and at $\snr = 15$~dB it sits $8\sigma$ below the baseline ($1.21 \times 10^{-2}$ versus $1.98 \times 10^{-2}$). Second, the learned head covers the majority of the baseline-to-oracle gap ($61\%/81\%/98\%$ at $\snr = 15/20/25$~dB), so the $\approx 2$~dB visible log-BER separation observed in Experiment~1 corresponds to a top-$\Ppath$ support-refinement module that is within a small fraction of the tightest support-oracle ceiling on the proposed dictionary. The residual gap between the learned head and the PCSI+MMSE reference is a detector-side gain-refinement and MP-BP residual that any support-refinement method on this dictionary cannot close.

\subsection{Experiment 4 --- complexity and MAC count}
\label{ssec:complexity_results}

Table~\ref{tab:complexity} reports the per-frame arithmetic cost of each receiver under two complementary lenses: the dominant-operator asymptotic complexity, and a first-order complex-MAC count at $\Nsub = 1024$. Three hypothetical recipes (re-OMP-per-iteration, dense LMMSE channel estimation, and the PCSI+MMSE oracle detector) are tabulated alongside the three deployed receivers to make the constant-factor differences that the asymptotic analysis hides explicit.

\begin{table}[t]
\centering
\caption{Complexity comparison for the receiver recipes considered in this report. ``Dominant operator'' gives the step that dominates the per-frame cost. MAC counts are first-order complex-multiplication estimates at $\Nsub = 1024$; runtime, memory, and latency effects on a realistic embedded target are out of scope for this report.}
\label{tab:complexity}
\small
\resizebox{\linewidth}{!}{%
\begin{tabular}{lllrr}
\toprule
Recipe & Dominant operator & Asymptotic complexity & MAC per frame & Relative \\
\midrule
\Serial{} (OMP + 1 MP-BP) & OMP & $\mathcal{O}(\Ppath \Nsub^2)$ & $\approx$ 14.0~M & $0.75\times$ \\
\ClassicalJoint{} (OMP-once + $5\times$ LS/MP-BP) & OMP + $5\times$ LS/MP-BP & $\mathcal{O}(\Ppath \Nsub^2)$ & $\approx$ 18.6~M & $\mathbf{1.00\times}$ (baseline) \\
\textbf{\LearnedHead{} (proposed)} & classical + $2\times$ head refine & $\mathcal{O}(\Ppath \Nsub^2)$ & $\approx$ 20.1~M & $\mathbf{\approx 1.08\times}$ \\
Re-OMP-per-iter (hypothetical) & $5\times$ OMP + $5\times$ LS/MP-BP & $\mathcal{O}(\Ppath \Nsub^2)$ ($5\times$ const.) & $\approx$ 48.6~M & $2.61\times$ \\
Dense LMMSE CE & $(\mbPhi^H\mbPhi)^{-1}$ & $\mathcal{O}(\Nsub^3)$ & $\approx$ 2.7~G & $\approx 145\times$ \\
\Oracle{} (dense MMSE, PCSI) & $(\Heff^H\Heff)^{-1}$ & $\mathcal{O}(\Nsub^3)$ & $\approx$ 2.7~G & $\approx 145\times$ (not deployable) \\
\bottomrule
\end{tabular}
}
\end{table}

The proposed learned receiver adds $\approx 1.5$~M MAC on top of the $\approx 18.6$~M-MAC classical core, an $\approx 8\%$ relative overhead, with the OMP step itself remaining the dominant per-frame cost in both recipes. Both dense alternatives ($\mathcal{O}(\Nsub^3)$) cost two orders of magnitude more per frame at $\Nsub = 1024$, while a hypothetical re-OMP-per-iteration schedule is a $5\times$ constant-factor more expensive than the classical joint without changing the asymptotic class. These constant-factor differences are the quantitative form of the claim that the learned support-refinement module delivers a non-trivial BER reduction without altering the sparse-recovery complexity profile.

The analytical MAC accounting of Table~\ref{tab:complexity} is complemented by a wall-clock measurement on three problem sizes ($\Nsub \in \{512, 1024, 2048\}$) reported in Table~\ref{tab:runtime}. Receiver invocation is timed on an NVIDIA RTX 4090 (single-GPU, complex64, $50$ frames per cell, channel-generation and data preparation excluded). The same configuration also benchmarks an externally-defined SP-AFDM receiver (Zheng \emph{et al.}~\cite{qing2024spafdm}) at its native operating point ($\Nsub = 512$, $\alpha_\text{max} = 2$, $\lmax = 1$, $Q + 1 = 10$, $M + 1 = 16$ pilots, $n_\text{iter} = 2$); the constraint $M(Q+1) < \Nsub - Q$ on Zheng's pilot scheme is incompatible with the present operating point's $Q \approx \Nsub / 3$ at $\Nsub \geq 1024$, so the SP-AFDM cell is informative as a cross-check on the absolute timing scale at $\Nsub = 512$ rather than as a head-to-head pivot. The learned head adds at most $31$~ms ($\approx 3.8\%$) over classical-Torch at $\Nsub = 2048$, and is at parity with classical-Torch at $\Nsub = 1024$ ($402$~ms vs.\ $416$~ms); the analytical $\approx 8\%$ MAC overhead is therefore not seen on RTX 4090 wall-clock at $\Nsub \leq 1024$ because GPU launch overhead and kernel fusion absorb the small MAC delta.

\begin{table}[t]
\centering
\caption{Wall-clock runtime per receiver invocation (mean, RTX 4090, complex64, $50$~frames per cell, channel generation and data preparation excluded). Peak GPU memory is reported for the $\Nsub = 2048$ cell. Source artefact: \texttt{outputs\_complexity/runtime\_table.\{csv,json,md\}}.}
\label{tab:runtime}
\small
\begin{tabular}{lrrrr}
\toprule
Receiver & $\Nsub = 512$ & $\Nsub = 1024$ & $\Nsub = 2048$ & Peak GPU mem ($\Nsub = 2048$) \\
\midrule
Classical NumPy joint OMP-MP (CPU) & $207$~ms & $463$~ms & $1\,855$~ms & --- \\
Classical Torch joint OMP-MP (GPU) & $242$~ms & $416$~ms & $\;\;\;803$~ms & $355$~MB \\
\textbf{Learned (C2-rel + F2-P, GPU)} & $257$~ms & $\mathbf{402}$~ms & $\;\;\;834$~ms & $368$~MB \\
Zheng \emph{et al.}~\cite{qing2024spafdm} SP-AFDM (native, GPU) & $268$~ms & --- & --- & --- \\
\bottomrule
\end{tabular}
\end{table}

\subsection{Summary}
\label{ssec:results_summary}

Three findings organise the quantitative results of this chapter. First, the one-shot serial receiver saturates at $\ber \approx 5 \times 10^{-2}$ from $\snr \geq 15$~dB with a frozen $\approx -11.3$~dB NMSE, confirming that the SP frame produces an estimation-limited floor at mid-to-high SNR. Second, the OMP-once + iterative LS/MP-BP recipe recovers approximately two to three orders of magnitude of BER relative to the serial receiver at $\snr = 25$~dB, while the learned support-refinement head adds a further consistent $\approx 2$~dB paired separation at $\snr \in \{20, 25, 30\}$~dB (3-seed pooled relative reductions $-37.7\%/-40.2\%/-37.6\%$, all statistically significant under continuity-corrected paired McNemar). Third, the two design analyses establish the mechanism: the single OMP pass recovers $\approx 95\%$ of the Dirichlet-expanded oracle support on average, with the missing $\approx 5\%$ structurally rather than noise-driven, and the learned head closes $61\%/81\%/98\%$ of the baseline-to-physical-oracle gap at $\snr = 15/20/25$~dB.
